% ATS-friendly cover letter - compile with pdfLaTeX (e.g. Overleaf)
\documentclass[10pt,a4paper]{article}
\usepackage[a4paper,left=1.9cm,right=1.9cm,top=2.0cm,bottom=1.6cm]{geometry}
\usepackage[T1]{fontenc}
\usepackage[utf8]{inputenc}
\usepackage{helvet}
\renewcommand{\familydefault}{\sfdefault}
\usepackage[hidelinks]{hyperref}
\input{glyphtounicode}
\pdfgentounicode=1

\pagestyle{empty}
\setlength{\parindent}{0pt}
\setlength{\parskip}{8pt}
\hypersetup{pdftitle={Arunabha Majumder - Cover Letter - Mechanical Design Engineer},pdfauthor={Arunabha Majumder}}

\begin{document}

{\LARGE\bfseries Arunabha Majumder}\\[2pt]
\textbf{Mechanical Design Engineer \textbar{} Mechatronic Systems \textbar{} Electromechanical Integration \textbar{} Prototyping \& Testing}\\[2pt]
{\small Aalborg, Denmark \textbar{} +45 71 62 24 00 \textbar{} \href{mailto:somrkmv1997@gmail.com}{somrkmv1997@gmail.com} \textbar{} \href{https://www.linkedin.com/in/arunabha-majumder-681264107}{linkedin.com/in/arunabha-majumder-681264107}\\
Portfolio: \href{https://arunabha22.github.io/arunabha-majumder/}{arunabha22.github.io/arunabha-majumder} \textbar{} Project: \href{https://www.viexo.aau.dk/}{www.viexo.aau.dk}}

September 29, 2026

Hiring Team, Engineering\\
GE Aerospace\\
United Kingdom

\textbf{Application: Mechanical Design Engineer -- Mechatronic Systems}

Dear Hiring Team,

I am applying for the Mechanical Design Engineer position with a focus on mechatronic systems. I hold a BTech in Mechanical Engineering and an MTech in Mechatronics, and I am finishing my PhD at Aalborg University. My work sits exactly where this role does: strong mechanical design, combined with sensing, actuation and control.

In my PhD, I designed a hybrid-actuated shoulder exoskeleton from concept to tested hardware. Rather than relying on a larger motor, I combined a parallel spring with a small 6\,Nm motor. I designed the components and assemblies in SolidWorks, used FEA for strength, stiffness and weight, made the parts through machining and additive manufacturing, and integrated the motor, encoder, force sensors, DAQ and embedded control. The design cut energy consumption by more than 25\% and won the Best Paper Award at IFToMM ISRM 2026.

I am comfortable working across the mechanical and electrical boundary. On a 5-bar planar parallel robot, I used variable-stiffness actuators to estimate the force at the end effector through the robot's Jacobian, and validated the estimate against a force sensor to within 10\% RMS error. I also designed the Assist-as-Needed controller for the exoskeleton, and during my Master's at CSIR-CMERI I built a pneumatic actuator test set-up and developed a controller that reduced tracking error to below 0.8\%.

I enjoy the workshop side of engineering. I build my own prototypes and test rigs, and I planned and ran a 12-participant validation campaign across several operating modes and load conditions. I have also taught and supervised Bachelor's students, and published peer-reviewed papers, which has sharpened how I communicate technical work to different audiences.

My CAD experience is in SolidWorks and Autodesk Fusion 360 rather than Inventor, and PCB design is an area I want to develop further. I am confident I can pick these up quickly, and I would value the chance to help build GE Aerospace's mechatronics capability while growing my own.

I would welcome the opportunity to discuss how I can contribute to your team. References are available upon request.

Kind regards,\\
Arunabha Majumder

\end{document}
